\chapter{Modèle vierge de dossier de sprint}
\begin{notebox}
\textbf{Mode d’emploi :} copier ce modèle dans un nouveau fichier \texttt{chapters/sprintXX.tex}, compléter les sections et produire les figures TikZ associées dans \texttt{figures/}. Les exemples S00 et S01 montrent la mise en page attendue.
\end{notebox}
\section{Fiche de sprint à remplir}
\begin{table}[htbp]\centering\small
\caption{Gabarit de fiche de sprint}
\begin{tabularx}{\textwidth}{@{}P{40mm}Y@{}}\toprule\textbf{Champ} & \textbf{À renseigner}\\\midrule
Sprint / version & Code, période et statut.\\
Objectif & Résultat utilisateur observable.\\
Stories incluses & Identifiants stables depuis le backlog.\\
Cas d’utilisation & UC-SXX-01, UC-SXX-02, etc.\\
Architecture & Classes, API, écrans et relations nécessaires.\\
Livrables & Code, maquettes, diagrammes, documentation.\\
Tests & Scénario nominal, erreurs, permissions, régressions.\\
Critère d’acceptation & Conditions démontrables, sans données fictives.\\
\bottomrule\end{tabularx}\end{table}

\section{Fiche détaillée de cas d’utilisation --- gabarit}
\begin{table}[htbp]\centering\small
\caption{Trame standard des cas d’utilisation}
\begin{tabularx}{\textwidth}{@{}P{36mm}Y@{}}\toprule\textbf{Champ} & \textbf{Description à compléter}\\\midrule
Code / intitulé & UC-SXX-YY : verbe et objet métier.\\
Acteur & Acteur principal et acteurs secondaires.\\
But & Valeur recherchée.\\
Préconditions & État requis avant exécution.\\
Déclencheur & Action initiatrice.\\
Scénario nominal & Séquence numérotée : actions acteur / réponses système.\\
Alternatives & Parcours valides de remplacement.\\
Exceptions & Erreurs, permissions et indisponibilités.\\
Postconditions & Effet sur l’état système.\\
Règles métier & Identifiants RG-XX et validations.\\
Critères de recette & Tests vérifiables, succès / échec.\\
\bottomrule\end{tabularx}\end{table}

\section{Livrables graphiques à joindre}
\begin{enumerate}
\item \textbf{Use case du sprint} : acteurs et fonctionnalités de l’incrément.
\item \textbf{Séquence de chaque cas détaillé important} : messages et réponses des composants.
\item \textbf{Diagramme de classes} : attributs, associations, multiplicités.
\item \textbf{Diagramme d’activité} : étapes, conditions et parcours d’erreur.
\item \textbf{Maquettes} : page d’accueil, liste, détail, formulaires, back-office ou autres écrans pertinents.
\end{enumerate}
\chapter{Sources et ressources de référence}
\begin{enumerate}
\item Documents de cadrage et backlog de \textbf{Souk El Hayawan} conservés dans Google Drive (documents de travail).
\item Tayara : \url{https://www.tayara.tn/}.
\item Sawt El Falleh : \url{https://www.sawtelfalleh.tn/}.
\item TunisieVet : \url{https://tunisievet.com/fr/vente-animaux}.
\item Animalia : \url{https://www.animalia.tn/}.
\item \textit{Farm Ghar}, rapport PFE fourni par l’utilisateur : référence de structure, non recopiée et non présentée comme source d’observations en Tunisie.
\item \textit{UML Specification}, Object Management Group : documentation de référence pour les diagrammes de conception.
\item Documentation officielle Spring Boot, Angular, PostgreSQL, Flyway et Keycloak : à citer avec versions validées lors de l’implémentation.
\end{enumerate}
